\documentclass[letterpaper,11pt]{article}

\usepackage[empty]{fullpage}
\usepackage{titlesec}
\usepackage{enumitem}
\usepackage[hidelinks]{hyperref}
\usepackage{fancyhdr}
\usepackage{fontawesome5}
\usepackage[T1]{fontenc}

\input{glyphtounicode}

\pagestyle{fancy}
\fancyhf{}
\renewcommand{\headrulewidth}{0pt}
\renewcommand{\footrulewidth}{0pt}

\addtolength{\oddsidemargin}{-0.5in}
\addtolength{\evensidemargin}{-0.5in}
\addtolength{\textwidth}{1in}
\addtolength{\topmargin}{-1in}
\addtolength{\textheight}{1.0in}
\setlength{\footskip}{4.08003pt}

\urlstyle{same}
\raggedbottom
\raggedright
\setlength{\tabcolsep}{0in}

\titleformat{\section}{
  \vspace{-6pt}\scshape\raggedright\large
}{}{0em}{}[\titlerule \vspace{-6pt}]

\pdfgentounicode=1

\newcommand{\resumeItem}[1]{\item\small{#1\vspace{-2pt}}}
\newcommand{\resumeSubheading}[4]{
  \vspace{-2pt}\item
    \begin{tabular*}{0.97\textwidth}[t]{l@{\extracolsep{\fill}}r}
      \textbf{#1} & #2 \\
      \textit{\small#3} & \textit{\small #4} \\
    \end{tabular*}\vspace{-7pt}
}
\newcommand{\resumeProjectHeading}[2]{
    \item
    \begin{tabular*}{0.97\textwidth}{l@{\extracolsep{\fill}}r}
      \small#1 & #2 \\
    \end{tabular*}\vspace{-7pt}
}
\newcommand{\resumeSubHeadingListStart}{\begin{itemize}[leftmargin=0.15in, label={}]}
\newcommand{\resumeSubHeadingListEnd}{\end{itemize}}
\newcommand{\resumeItemListStart}{\begin{itemize}}
\newcommand{\resumeItemListEnd}{\end{itemize}\vspace{-5pt}}
\renewcommand\labelitemii{$\vcenter{\hbox{\tiny$\bullet$}}$}

\begin{document}

\vspace*{12pt}
\begin{center}
    \textbf{\Huge \scshape Alireza Noroozi \\ \vspace{8pt}}
    \faIcon{envelope}
    \href{mailto:noroozi.alirezaa@gmail.com}{\underline{noroozi.alirezaa@gmail.com}}
    \faIcon{github}
    \href{https://github.com/Alirzeanoroozi}{\underline{Github}}
    \faIcon{blog}
    \href{https://alirzeanoroozi.github.io/}{\underline{Personal Page}}
    \faIcon{google}
    \href{https://scholar.google.com/citations?user=CzouC0IAAAAJ\&hl=en}{\underline{Google Scholar}}
\end{center}

\section{Research Interests}
    \resumeItemListStart
        \resumeItem{Peptide and protein binder design; generative models for biological sequences; discrete and insertion--deletion diffusion; evaluation of designed binders with molecular dynamics; protein--ligand cofolding at adaptive resolution.}
    \resumeItemListEnd

\section{Education}
  \resumeSubHeadingListStart
    \resumeSubheading
      {Ko\c{c} University}{October 2025 -- June 2027 (expected)}
      {MSc in Computational Sciences and Engineering}{GPA: 3.48/4.0}
    \resumeSubheading
      {Sharif University of Technology}{September 2020 -- June 2024}
      {Bachelor of Science in Computer Engineering}{GPA: 17.32/20 (3.46/4.0)}
  \resumeSubHeadingListEnd

\section{Publications and Preprints}
    \resumeItemListStart
        \resumeItem{\textbf{A. Noroozi}, \"{O}. Keskin, A. G\"{u}rsoy, Y. Karami, and H. Khakzad, \textbf{IDPepDesign: Insertion-Deletion is All Peptide Design Needs}, preprint, 2026; ML4SM Workshop, NeurIPS 2026.}
        \resumeItem{\textbf{BinderEval}: evaluating designed protein binders with molecular dynamics data and COMPASS, with the DELTA team, Inria Nancy, preprint, 2026.}
        \resumeItem{M. Ghasemiagreh and \textbf{A. Noroozi}, \textbf{BoltzGen-ADC}: Generative Design of Universal Protease-Cleavable Peptide Linkers for Antibody--Drug Conjugates, Merck Challenge submission, July 2026.}
        \resumeItem{\textbf{A. Noroozi}, \"{O}. Keskin, and A. G\"{u}rsoy, \textbf{AMP-DiT}: Antimicrobial Peptide Design with AMP-classifier Conditional Diffusion Transformers, GenBio Workshop at ICML 2026 (poster). \href{https://openreview.net/forum?id=KYzd6lcfZ0}{\underline{paper}} \href{https://github.com/Alirzeanoroozi/DiT-AMP}{\underline{code}}}
        \resumeItem{M. Kalemati, \textbf{A. Noroozi}, A. Shahbakhsh et al., ParaAntiProt provides paratope prediction using antibody and protein language models, \emph{Scientific Reports} 14, 29141 (2024). \href{https://doi.org/10.1038/s41598-024-80940-y}{\underline{link}}}
        \resumeItem{J. Reyes, \textbf{A. Noroozi}, Y. Xiao, and M. Kersten-Oertel, Bridging the Gaps: Imputation of Parkinson's Disease Clinical Assessments With Federated Learning, \emph{IEEE Journal of Biomedical and Health Informatics} (2025). \href{https://ieeexplore.ieee.org/document/11098925}{\underline{link}}}
    \resumeItemListEnd

\section{Research Experience}
    \resumeSubHeadingListStart
        \resumeProjectHeading
        {\textbf{DELTA Team, Inria Nancy} $|$ \footnotesize\emph{Peptide Binder Design, Molecular Dynamics}}{June 2026 -- Present}
        \resumeItemListStart
            \resumeItem{Research internship under the supervision of Hamed Khakzad and Yasaman Karami}
            \resumeItem{Developed \textbf{IDPepDesign}, a target-conditioned insertion--deletion diffusion model for peptide binder design, built on SCISOR with a pooled target representation injected through denoising and residue-level cross-attention; improved perplexity 17.1 $\rightarrow$ 11.3 and Fr\'{e}chet protein distance 1.84 $\rightarrow$ 0.75, with competitive predicted interface quality across 11 benchmark targets}
            \resumeItem{Contributed to \textbf{BinderEval}, which evaluates designed binders using molecular dynamics data and COMPASS instead of single-structure confidence scores, and released a new version of \textbf{COMPASS}}
            \resumeItem{Invited talk at PhD Pizza Time, Inria Nancy (September 2026)}
        \resumeItemListEnd

        \resumeProjectHeading
        {\textbf{AMP-DiT, Ko\c{c} University} $|$ \footnotesize\emph{Discrete Diffusion, Generative Modeling}}{October 2025 -- Present}
        \resumeItemListStart
            \resumeItem{Computational Sciences and Engineering; supervised by Prof. Attila G\"{u}rsoy and Prof. \"{O}zlem Keskin}
            \resumeItem{Built a conditional discrete diffusion framework that generates antimicrobial peptides directly in sequence space, avoiding the representation mismatch of protein language models trained on full-length proteins}
            \resumeItem{Used classifier conditioning from Macrel-derived activity and non-hemolytic safety scores; conditioning on a single predictor generalised across multiple AMP predictors, indicating learned antimicrobial features rather than scorer overfitting}
            \resumeItem{Accepted as a poster at the GenBio Workshop at ICML 2026; submitted to the AMP Challenge}
        \resumeItemListEnd

        \resumeProjectHeading
        {\textbf{AMP-BindGen, Ko\c{c} University} $|$ \footnotesize\emph{Targeted Binder Design, AMR}}{2026 -- Present}
        \resumeItemListStart
            \resumeItem{Guided design of peptide-like binders and inhibitors against the clinically important carbapenemases \textbf{NDM-5} and \textbf{KPC-3}, in collaboration with KUISCID and the Ko\c{c} University School of Medicine}
            \resumeItem{Designed a composite ranking scheme over interface confidence, catalytic-site engagement, predicted antimicrobial activity and toxicity filters for candidate selection}
            \resumeItem{Designed peptides have been validated experimentally in the lab, alongside metrics for peptide--protein interaction quality}
            \resumeItem{Presented as posters at ECCB 2026 (Geneva) and the Ko\c{c} University AI Symposium}
        \resumeItemListEnd

        \resumeProjectHeading
        {\textbf{TED, Ko\c{c} University / 310.ai} $|$ \footnotesize\emph{Protein Domain Prediction, CATH Classification}}{2024 -- Present}
        \resumeItemListStart
            \resumeItem{\href{https://github.com/Alirzeanoroozi/TED}{\underline{Project link}}}
            \resumeItem{End-to-end sequence-to-sequence transformer that predicts domain boundaries and CATH identifiers directly from sequence, removing the external search tools other methods depend on}
            \resumeItem{Trained on $\sim$186M unique protein sequences from the TED database with LoRA fine-tuning over ESM-2 embeddings; ran large-scale preprocessing and fine-tuning on GPU clusters}
        \resumeItemListEnd

        \resumeProjectHeading
        {\textbf{BoltzGen-ADC} $|$ \footnotesize\emph{Antibody--Drug Conjugates, Generative Design}}{July 2026}
        \resumeItemListStart
            \resumeItem{Pipeline over BoltzGen and ADFLIP that designs short protease-cleavable peptide linkers for an arbitrary antibody variable region and small-molecule payload given as isomeric SMILES, without retraining}
            \resumeItem{Added filters for Fc preservation, protease-cleavage liability, and Boltz-based validation of peptide--protease complexes; demonstrated on brentuximab--MMAE conjugates}
        \resumeItemListEnd

        \resumeProjectHeading
        {\textbf{Taxonomy Lineage Generation, 310.ai} $|$ \footnotesize\emph{Transformers, Protein Language Models}}{2024 -- 2025}
        \resumeItemListStart
            \resumeItem{Built a transformer that generates full taxonomy lineages from protein sequences}
            \resumeItem{Designed training pipelines, tokenisation strategies, and loss functions to stabilise sequence generation}
            \resumeItem{Addressed token collapse and improved generalisation with curriculum learning and embedding regularisation}
        \resumeItemListEnd

        \resumeProjectHeading
        {\textbf{IgG1 Binder Design, 310.ai} $|$ \footnotesize\emph{Protein Design, Structural Modeling}}{2024 -- 2025}
        \resumeItemListStart
            \resumeItem{Designed IgG1 antibody binders using generative protein models and sequence design pipelines}
            \resumeItem{Integrated structural constraints and protein language model embeddings to guide binder specificity and stability}
        \resumeItemListEnd

        \resumeProjectHeading
        {\textbf{OPL, Sharif University} $|$ \footnotesize\emph{Transfer Learning, Antibody, Paratope Prediction}}{June 2023 -- September 2024}
        \resumeItemListStart
            \resumeItem{Department of Computer Engineering; bachelor project under the supervision of Dr. Koohi}
            \resumeItem{\textbf{ParaAntiProt}: transfer learning for paratope prediction from sequence alone}
            \resumeItem{Improved paratope prediction accuracy using antibody and protein language models; published in \emph{Scientific Reports}}
        \resumeItemListEnd

        \resumeProjectHeading
        {\textbf{apLab, Concordia University} $|$ \footnotesize\emph{Federated Learning, Medical ML}}{June 2023 -- September 2023}
        \resumeItemListStart
            \resumeItem{Department of Computer Science, Montreal, Canada; internship under the supervision of Dr. Kersten-Oertel}
            \resumeItem{Compared and implemented federated learning algorithms for data imputation on Parkinson's disease clinical assessments; published in \emph{IEEE JBHI}}
            \resumeItem{Implemented PW, FedAvg, and FedOpt in PyTorch, including work on a convergence proof for PW}
            \resumeItem{Developed federated algorithms for non-identically distributed data, and used GANs for prediction from distributed data}
        \resumeItemListEnd

        \resumeProjectHeading
        {\textbf{OPL, Sharif University} $|$ \footnotesize\emph{Attention Models, MHC--Peptide Binding Affinity}}{December 2022 -- March 2023}
        \resumeItemListStart
            \resumeItem{Department of Computer Engineering; bachelor student under the supervision of Dr. Koohi}
            \resumeItem{MHC--peptide binding prediction using attention mechanisms}
        \resumeItemListEnd
    \resumeSubHeadingListEnd

\section{Talks and Posters}
    \resumeItemListStart
        \resumeItem{Talk, \textbf{PhD Pizza Time}, Inria Nancy, September 2026.}
        \resumeItem{Posters, \textbf{ECCB 2026}, Geneva --- AMP-BindGen and AMP-DiT, September 2026.}
        \resumeItem{Posters, \textbf{Ko\c{c} University AI Symposium}, April 2026 --- AMP-BindGen and TED; joint talk on treatment options for carbapenem resistance and AMP design with AI.}
        \resumeItem{Poster, \textbf{GenBio Workshop at ICML 2026} --- AMP-DiT.}
    \resumeItemListEnd

\section{Working Experience}
    \resumeSubHeadingListStart
        \resumeProjectHeading
        {\textbf{310.ai} $|$ \footnotesize\emph{AI Engineer in Computational Biology}}{September 2024 -- October 2025}
        \resumeItemListStart
            \resumeItem{Deep learning models for protein design, including taxonomy prediction and binder design pipelines}
            \resumeItem{Developed an AI chatbot backend using LLM tool calling}
        \resumeItemListEnd

        \resumeProjectHeading
        {\textbf{Ershan Total Solution} $|$ \footnotesize\emph{PHP, Back-end, Web Developer}}{June 2022 -- September 2022}
        \resumeItemListStart
            \resumeItem{Developed a stock trading website using the Laminas framework in PHP}
            \resumeItem{Built an Excel and PDF generator for all reports}
        \resumeItemListEnd

        \resumeProjectHeading
        {\textbf{AI Medical} $|$ \footnotesize\emph{Deep Learning, Medical Imaging}}{June 2021 -- September 2021}
        \resumeItemListStart
            \resumeItem{Summer internship in machine learning for medical data and medical imaging}
        \resumeItemListEnd

        \resumeProjectHeading
        {\textbf{Salam High School} $|$ \footnotesize\emph{Mathematics Olympiad Tutor}}{September 2021 -- March 2023}
        \resumeItemListStart
            \resumeItem{Taught mathematical olympiad topics (geometry, algebra, combinatorics, number theory)}
        \resumeItemListEnd
    \resumeSubHeadingListEnd

\section{Honors}
    \resumeItemListStart
        \resumeItem{Awarded the 2025 Presidential Fellowship at Ko\c{c} University for MSc entrance.}
        \resumeItem{Awarded the 2024 Vector Scholarship for MSc entrance at York University, valued at 17,500 CAD (deferred).}
        \resumeItem{Ranked 31st among 150k+ participants in Iran's national university entrance exam in 2020 (Konkur).}
        \resumeItem{Silver (2019) and Bronze (2018) medals in \textbf{Iran's National Mathematical Olympiad}.}
    \resumeItemListEnd

\section{Reviewing}
\resumeItemListStart
    \resumeItem{Reviewer, GenBio Workshop (Generative and Agentic AI for Biology), ICML 2026}
    \resumeItem{Reviewer, ML4SM Workshop, NeurIPS 2026}
\resumeItemListEnd

\section{Teaching Assistant Experience}
    \resumeItemListStart
        \resumeItem{Teaching Assistant, Python Programming, Ko\c{c} University}
        \resumeItem{Head of Teaching Assistant Team, Probability and Statistics, 2 semesters (lecturer: Dr. Sharifi-Zarchi)}
        \resumeItem{Head of Teaching Assistant Team, Compiler Design (lecturer: Dr. Hosseinmardi)}
        \resumeItem{Projects Head, Teaching Assistant Team, Introduction to Bioinformatics}
        \resumeItem{Teaching Assistant, Discrete Structures and Basic Programming in Python}
    \resumeItemListEnd

\section{Technical Skills}
    \resumeItemListStart
        \resumeItem{\textbf{Languages:} Python, C++, JavaScript, PHP}
        \resumeItem{\textbf{Machine learning:} PyTorch, Hugging Face, scikit-learn, distributed and multi-GPU training, SLURM clusters}
        \resumeItem{\textbf{Structural biology:} Boltz, BoltzGen, ADFLIP, AlphaFold-class predictors, RFdiffusion, ESM-2/ESM-C, ProtT5, Biopython, molecular dynamics analysis}
        \resumeItem{\textbf{Other:} Git, Docker, Linux}
    \resumeItemListEnd

\section{Language Proficiency}
    \resumeItemListStart
        \resumeItem{Persian: Native}
        \resumeItem{English: Fluent, TOEFL (T: 97, R: 26, L: 25, S: 25, W: 21)}
    \resumeItemListEnd

\end{document}
